\subsection{准素分解的局部化}

给定诺特环 A 上模 M 的一个子模 N，为简便起见，本节将 \(D_{I}(M/N)\) 记为 N 在 M 中的、索引集为 I（与 Ass(M/N) 等势）的既约准素分解所成的集合。

命题 6. 设 A 是诺特环，M 是 A-模，N 是 M 的子模，且 I = Ass(M/N)。设 S 是 A 的乘性子集，J 是 I 的子集，由指标 \(i\) 组成，且这些指标满足 \(\mathbf{S} \cap \mathfrak{p}_i = \varnothing\)。令 \(\mathbf{N}'\) 是 \(\mathbf{N}\) 关于 \(\mathbf{S}\) 在 \(\mathbf{M}\) 中的饱和。则：

\begin{enumerate}
\def\labelenumi{(\roman{enumi})}
\item
  如果 \((\mathbf{Q}_i)_{i\in I}\) 是 \(\mathbf{D}_{\mathrm{I}}(\mathbf{M} / \mathbf{N})\) 的元素，则族 \((\mathbf{Q}_i)_{i\in J}\) 是 \(\mathbf{D}_{J}(\mathbf{M} / \mathbf{N}')\) 的元素，且族 \((\mathbf{S}^{-1}\mathbf{Q}_i)_{i\in J}\) 是 \(\mathbf{D}_{J}(\mathbf{S}^{-1}\mathbf{M} / \mathbf{S}^{-1}\mathbf{N})\) 的元素。
\item
  映射 \((\mathbf{Q}_i)_{i\in J}\to (\mathbf{S}^{-1}\mathbf{Q}_i)_{i\in J}\) 是从 \(\mathrm{D}_{\mathrm{J}}(\mathbf{M} / \mathbf{N}^{\prime})\) 到 \(\mathbf{D}_{\mathrm{J}}(\mathbf{S}^{-1}\mathbf{M} / \mathbf{S}^{-1}\mathbf{N})\) 的双射。
\item
   如果\((\mathbf{Q}_i)_{i\in J}\)是\(\mathbf{D}_{\mathbf{J}}(\mathbf{M} / \mathbf{N}')\)的元素，而\((\mathbf{R}_i)_{i\in I}\)是\(\mathbf{D}_{\mathbf{I}}(\mathbf{M} / \mathbf{N})\)的元素，族\((\mathrm{T}_i)_{i\in I}\)满足\(\mathrm{T}_i = \mathrm{Q}_i\)对于\(i\in \mathbf{J}\)，且\(\mathrm{T}_i = \mathrm{R}_i\)对于\(i\in \mathbf{I} - \mathbf{J}\)的元素\(\mathbf{D}_{\mathbf{I}}(\mathbf{M} / \mathbf{N})\)。
\item
   我们知道（第 1 节，命题 3），对于\(i \in J\)，\(S^{-1}Q_{i}\)是关于\(S^{-1}p_{i}\)的准素子模，并且对于\(i \in I - J\)，\(S^{-1}Q_{i} = S^{-1}M\)；由于\(S^{-1}N = \bigcap_{i \in I} S^{-1}Q_{i}\)（第二章，第 2 节，第 4 节），所以也有\(S^{-1}N = \bigcap_{i \in J} S^{-1}Q_{i}\)。对于\(S^{-1}p_{i}\)，\(i \in J\)彼此不同，且其集合是\(\operatorname{Ass}(S^{-1}M / S^{-1}N)\)（§ 1，第 2 节，命题 5 的推论）；于是（命题 4），\(S^{-1}Q_{i})_{i \in J}\)是\(S^{-1}N\)的既约准素分解。此外，\(Q_{i} = (i_{M}^{S})^{-1}(S^{-1}Q_{i})\)（第 1 节命题 3），所以\(N' = (i_{M}^{S})^{-1}(S^{-1}N) = \bigcap_{i \in J} Q_{i}\)，而\((Q_{i})_{i \in J}\)显然是\(N'\)中的既约准素分解。
\item
   由于\(S^{-1}N' = S^{-1}N\)，我们可以用\(N'\)替换 N，即设 J = I。令\((P_i)_{i \in I}\)是\(S^{-1}N\)在\(S^{-1}M\)中的既约准素分解，并记\(Q_i = (i_M^S)^{-1}(P_i)\)；根据第 1 节命题 3，\(Q_i\)是关于\(p_i (i \in I)\)的准素子模，并且\((Q_i)_{i \in I}\)随后是 N 在 M 中的既约准素分解。最后，对所有\(i \in I\)以及 M 的每个子模\(Q'_i\)，若 M 关于\(p_i\)-准素，则\(Q'_i = (i_M^S)^{-1}(S^{-1}Q'_i)\)（根据第 1 节命题 3），我们定义了两个映射\(D_I(M/N) \to D_I(S^{-1}M/S^{-1}N)\)和\(D_I(S^{-1}M/S^{-1}N) \to D_I(M/N)\)，它们的复合分别是\(D_I(M/N)\)和\(D_I(S^{-1}M/S^{-1}N)\)上的恒等映射，这就证明了 (ii)。
\item
  根据 (i)，\(\mathbf{N}' = \bigcap_{i\in J}\mathbb{R}_i\)，从而
\end{enumerate}

\[
\mathrm{N} = \mathrm{N} ^ {\prime} \cap \bigcap_ {i \in \mathrm{I} - J} \mathrm{R} _ {i} = \left(\bigcap_ {i \in J} \mathrm{Q} _ {i}\right) \cap \left(\bigcap_ {i \in \mathrm{I} - J} \mathrm{R} _ {i}\right)
\]

并且由第 3 节命题 4 的推论立即可知该准素分解是既约的。

推论。映射

\[
\mathrm{D} _ {\mathrm{I}} (\mathrm{M} / \mathrm{N}) \rightarrow \mathrm{D} _ {\mathrm{J}} (\mathrm{M} / \mathrm{N} ^ {\prime}) \quad \text{且} \quad \mathrm{D} _ {\mathrm{I}} (\mathrm{M} / \mathrm{N}) \rightarrow \mathrm{D} _ {\mathrm{J}} (\mathrm{S} ^ {- 1} \mathrm{M} / \mathrm{S} ^ {- 1} \mathrm{N})
\]

在命题 6 (i) 中定义的映射是满射。

命题 6 (iii) 表明映射 \(\mathrm{D}_{\mathrm{I}}(\mathrm{M}/\mathrm{N}) \rightarrow \mathrm{D}_{\mathrm{J}}(\mathrm{M}/\mathrm{N}')\) 是满射，而命题 6 (ii) 随即表明映射

\[
\mathrm{D} _ {\mathrm{I}} (\mathrm{M} / \mathrm{N}) \rightarrow \mathrm{D} _ {\mathrm{J}} (\mathrm{S} ^ {- 1} \mathrm{M} / \mathrm{S} ^ {- 1} \mathrm{N})
\]

是满射。

